\section{Growth-controlled and local-to-global regularization}\label{sec:wu-discharge}

This chapter proves the two statements of \cite[Lemma~5.11(5),(6)]{FHX}
(Theorems~\ref{thm:import-growth-control} and~\ref{thm:import-local-to-global}) in three layers.
\begin{enumerate}
\item \emph{Affine reduction} (\S\ref{subsec:wu-affine}). Both statements follow, by arguments on
$Y^\an$ alone, from the approximation of convex PL functions that agree with a Fubini--Study function
outside a compact set (Theorem~\ref{thm:wu-pl-approx}). Conversely that statement is the special case of
\cite[Lemma~5.11(5),(6)]{FHX} for such functions (Lemma~\ref{lem:wu-pl-special-case}). No
compactification, no decreasing-limit theorem (PSH2), no countable regularization
\cite[Corollary~12.18]{BJ22} and no countability of $k$ are used.
\item \emph{Compactification} (\S\ref{subsec:wu-compactification}).
The uniform approximation statement of \cite[Lemma~3.24(1)]{Wu26}
follows from positivity on a normal projective compactification
built from the graded integral closure of a Rees-type algebra \cite[Lemma~5.9]{FHX}. Only continuous
functions that are Fubini--Study near the boundary are fed to positivity, so \cite[Lemma~5.10]{FHX} is
not used.
\item \emph{Positivity} (\S\ref{subsec:wu-positivity}). Positivity is \cite[Lemma~5.5]{FHX} combined
with the uniform approximation of continuous psh functions in an ample class \cite{BJ22}. The proof of
\cite[Lemma~5.5]{FHX} uses test configurations \cite{BJ22} and the positivity criterion of Chambert-Loir--Ducros
\cite{CLD12}; these are recorded below (Theorems~\ref{thm:wu-bj-testconfig}, \ref{thm:wu-bj-nef-psh}
and~\ref{thm:wu-cld-nef}) for comparison only. Here the lemma is proved through the native pipeline
(Theorem~\ref{thm:lpn-local-positivity-instance}), where nefness on the central fibre is proved without
\cite{CLD12}, and none of these three statements is used.
\end{enumerate}
No result of this chapter depends on Theorems~\ref{thm:import-growth-control}
and~\ref{thm:import-local-to-global}, except Lemma~\ref{lem:wu-pl-special-case}, which compares them with
Theorem~\ref{thm:wu-pl-approx}. Nothing depends on \cite[Theorem~6.21]{BJ22}.

\subsection{Affine reduction}\label{subsec:wu-affine}

\begin{convention}\label{conv:wu-setting}
\notready
\uses{conv:sec5, def:fs, lem:fs-basic, lem:polydisc-compact}
In this section $k$ is algebraically closed with the trivial absolute value; countability is not assumed.
$Y=\Spec A$ is a normal affine variety. A sequence $(u_j)_{j\in\N}$ of functions
$Y^\an\to[-\infty,\infty)$ is \emph{decreasing} if $u_{j+1}\le u_j$ pointwise; its limit is $\inf_ju_j$.
Every $g=\max\{a_1\log|f_1|+b_1,\dots,b_0\}\in\FS(Y)$ satisfies $g\ge b_0$, and $g-b_0$, $cg$ and $g+q$ lie in
$\FS(Y)$ for $c\in\Q_{>0}$, $q\in\Q$ (Lemma~\ref{lem:fs-basic}). Sublevel sets $\{g\le t\}$, $t\in\R$, are
compact: they are closed, and since $g\ge a_i\log|f_i|+b_i$ for the generators $f_i$ occurring in $g$
they lie in a polydisc (Lemma~\ref{lem:polydisc-compact}).
\end{convention}

\begin{definition}[Functions standard at infinity {\cite[Lemma~3.24]{Wu26}}]\label{def:wu-standard}
\notready
\uses{conv:wu-setting, def:convex-pl, def:fs}
Let $h\in\FS(Y)$. A function $w\colon Y^\an\to\R$ is \emph{standard at infinity with respect to $h$} if
$w$ is convex PL on $Y^\an$ and there is a compact set $E\subseteq Y^\an$ with $w=h$ on $Y^\an\setminus E$.
Such a $w$ is continuous, hence bounded on $E$; since $h\ge b_0$, there is $C\in\R$ with $w\le h+C$ on
$Y^\an$.
\end{definition}

\begin{theorem}[Approximation of convex PL functions standard at infinity {\cite[Lemma~5.11]{FHX}, \cite[Lemma~3.24]{Wu26}}]\label{thm:wu-pl-approx}
\lean{EnvelopeProperty.ConvexPLApproximation, EnvelopeProperty.LocalPositivity.exists_antitone_fsData, EnvelopeProperty.LocalPositivity.toConvexPLApproximation}
\leanok
\uses{def:wu-standard, thm:wu-pl-approx-uniform, lem:fs-basic}
Let $w$ be standard at infinity with respect to $h\in\FS(Y)$. For every rational $c>1$ there are
$h_j\in\FS(Y)$ decreasing to $w$ and constants $C_j\in\R$ with $h_j\le ch+C_j$.
\end{theorem}

\begin{proof}
Let $E$ be as in Definition~\ref{def:wu-standard} and $b_0$ the constant term of $h$. By
Theorem~\ref{thm:wu-pl-approx-uniform} choose $v_j\in\FS(Y)$ with $|v_j-w|\le2^{-j}$ on $Y^\an$, and put
$h_j\coloneqq v_j+3\cdot2^{-j}\in\FS(Y)$. Then
$h_{j+1}\le w+2^{-j-1}+3\cdot2^{-j-1}=w+2^{1-j}\le v_j+2^{-j}+2^{1-j}=h_j$ and $w\le h_j\le w+2^{2-j}$, so
the $h_j$ decrease to $w$. Off $E$ we have $h_j\le h+4\le ch-(c-1)b_0+4$, because $(c-1)h\ge(c-1)b_0$. On
$E$ we have $h_j\le\max_Ew+4$ and $ch\ge cb_0$. Hence $h_j\le ch+C_j$ with
$C_j\coloneqq\max\{4-(c-1)b_0,\ \max_Ew+4-cb_0\}$.
\end{proof}

\begin{lemma}[Comparison with the statements of Lemma~5.11(5),(6) {\cite[Lemma~5.11]{FHX}}]\label{lem:wu-pl-special-case}
\lean{EnvelopeProperty.ConvexPLApproximation.of_wuRegularization, EnvelopeProperty.exists_le_eval_add_of_isConvexPLOn_univ}
\leanok
\uses{def:wu-standard, thm:import-growth-control, thm:import-local-to-global, def:convex-pl}
The statements of Theorems~\ref{thm:import-growth-control} and~\ref{thm:import-local-to-global} imply the
conclusion of Theorem~\ref{thm:wu-pl-approx}.
\end{lemma}

\begin{proof}
Let $w$ be standard at infinity with respect to $h$. It is real valued, hence not identically $-\infty$; it
satisfies $w\le h+C$ (Definition~\ref{def:wu-standard}), so it has logarithmic growth; and on every open
set it is the limit of the constant sequence $w$, which is convex PL. Theorem~\ref{thm:import-local-to-global}
gives $w\in\Rcal(Y)$, and Theorem~\ref{thm:import-growth-control}, applied with $g=h$, gives for every
rational $c>1$ functions $u_j\in\FS(Y)$ decreasing to $w$ with $u_j-ch$ bounded above.
\end{proof}

\begin{lemma}[Collar gluing {\cite[Lemma~5.10]{FHX}, \cite[Lemma~3.25]{Wu26}}]\label{lem:wu-collar}
\lean{EnvelopeProperty.exists_collar}
\leanok
\uses{conv:wu-setting, def:wu-standard, lem:local-basic, lem:fs-basic}
Let $g\in\FS(Y)$ with $g\ge0$, let $c\in\Q$ with $c>1$, $C\in\R$, $C_1\in\Q$, and let
$v\colon Y^\an\to[-\infty,\infty)$ satisfy $v\le g+C$. Let $t\in\R$ with $(c-1)t\ge C-C_1+2$, let
$U\subseteq Y^\an$ be an open set containing $\{g\le t+1\}$, and let $(v_j)$ be a decreasing sequence of
convex PL functions on $U$ with limit $v$ on $U$. Put $U'\coloneqq\{g<t+1\}$. There is $j_0$ such that for
$j\ge j_0$ the function $W_j$ equal to $\max\{v_j,cg+C_1\}$ on $U'$ and to $cg+C_1$ on $Y^\an\setminus U'$
is standard at infinity with respect to $cg+C_1\in\FS(Y)$, with $W_j=cg+C_1$ on $\{g>t\}$; moreover
$(W_j)_{j\ge j_0}$ decreases to $\max\{v,cg+C_1\}$.
\end{lemma}

\begin{proof}
On $\{g\ge t\}$ we have $(c-1)g\ge C-C_1+2$, hence $v\le g+C\le cg+C_1-2$. The collar
$Q\coloneqq\{t\le g\le t+1\}$ is compact (Convention~\ref{conv:wu-setting}) and contained in $U$. The
functions $v_j$ are continuous on $U$ and decrease to $v<cg+C_1-1$ on $Q$; the sets
$\{y\in U:v_j(y)<cg(y)+C_1-1\}$ are open, increase with $j$ and cover $Q$, so one of them, say for $j_0$,
contains $Q$, and then so do those for $j\ge j_0$. Fix $j\ge j_0$. The set
$\{y\in U':v_j(y)>cg(y)+C_1\}$ is contained in $\{g<t\}$, so its closure is contained in $\{g\le t\}\subseteq U'$.
By the gluing statement of Lemma~\ref{lem:local-basic}, applied to the convex PL function $cg+C_1$ on $Y^\an$
and to $v_j$ on the open set $U'\subseteq U$, the function $W_j$ is convex PL on $Y^\an$. If $g(y)>t$ then
either $y\notin U'$, or $y\in Q$ and $v_j(y)<cg(y)+C_1$; in both cases $W_j(y)=cg(y)+C_1$. So $W_j$ is standard
at infinity with respect to $cg+C_1$, with $E=\{g\le t\}$. The $W_j$ decrease because the $v_j$ do. Their limit
is $\max\{v,cg+C_1\}$ on $U'$, and $cg+C_1=\max\{v,cg+C_1\}$ on $\{g\ge t+1\}$ since $v\le cg+C_1-2$ there.
\end{proof}

\begin{lemma}[Growth-controlled diagonal sequences]\label{lem:wu-growth-diagonal}
\lean{EnvelopeProperty.exists_antitone_fs_of_antitone_growth, EnvelopeProperty.exists_antitone_seq_of_directed, EnvelopeProperty.exists_lt_on_compact_of_antitone_of_continuousOn}
\leanok
\uses{conv:wu-setting, lem:fs-basic}
Let $g\in\FS(Y)$ with $g\ge0$ and let $0<c'<c$ be rationals. Let $(w_n)_{n\in\N}$ be a decreasing sequence of
functions $Y^\an\to[-\infty,\infty)$ with limit $w$, and suppose that each $w_n$ is the limit of a decreasing
sequence $(F_{n,j})_{j\in\N}$ in $\FS(Y)$ with $F_{n,0}\le c'g+D_n$ for some $D_n\in\R$. Then $w$ is the limit
of a decreasing sequence $(G_m)$ in $\FS(Y)$ such that $cg-A_m\le G_m\le cg+B_m$ for constants $A_m,B_m\in\R$.
\end{lemma}

\begin{proof}
For $m\in\N$ and all $n,j$ put $H^{(m)}_{n,j}\coloneqq\max\{F_{n,j},cg-m\}\in\FS(Y)$. Since
$F_{n,j}\le F_{n,0}\le c'g+D_n$, we have $H^{(m)}_{n,j}=cg-m$ outside the compact set
$K_m\coloneqq\{g\le(m+D_m)/(c-c')\}$ when $n=m$. We choose indices $j_m$ inductively and put
$G_m\coloneqq H^{(m)}_{m,j_m}+2^{-m}$, subject to
\begin{enumerate}
\item[(a)] $G_m\le G_{m-1}$ if $m\ge1$, and
\item[(b)] $G_m\le H^{(m)}_{l,m}+2^{1-m}$ for all $l\le m$.
\end{enumerate}
Off $K_m$, $G_m=cg-m+2^{-m}$, while $G_{m-1}\ge cg-m+1+2^{1-m}$ and $H^{(m)}_{l,m}+2^{1-m}\ge cg-m+2^{1-m}$
everywhere; so (a) and (b) hold off $K_m$ for every choice of $j_m$. On $K_m$, as $j\to\infty$ the continuous
real-valued functions $H^{(m)}_{m,j}+2^{-m}$ decrease to $\max\{w_m,cg-m\}+2^{-m}$, which is strictly smaller
than $G_{m-1}=\max\{F_{m-1,j_{m-1}},cg-m+1\}+2^{1-m}$ (as $w_m\le w_{m-1}\le F_{m-1,j_{m-1}}$) and than
$H^{(m)}_{l,m}+2^{1-m}$ for $l\le m$ (as $w_m\le w_l\le F_{l,m}$). These finitely many strict inequalities
between continuous functions hold on the compact set $K_m$ for all large $j$: the open sets where they hold
increase with $j$ and cover $K_m$. Choose such a $j=j_m$.

The $G_m$ lie in $\FS(Y)$ and decrease by (a). Since $G_m\ge\max\{w_m,cg-m\}\ge w$, the limit is $\ge w$.
For fixed $l$ and $x\in Y^\an$, (b) gives $\inf_mG_m(x)\le\inf_{m\ge l}\bigl(\max\{F_{l,m}(x),cg(x)-m\}+2^{1-m}\bigr)
=w_l(x)$, also when $w_l(x)=-\infty$; hence the limit is $\le\inf_lw_l=w$. Finally $G_m\ge cg-m$, and
$G_m-cg$ equals $-m+2^{-m}$ off $K_m$ and is continuous on the compact set $K_m$, hence bounded above.
\end{proof}

\begin{proposition}[Regularization with growth control {\cite[Lemma~5.11]{FHX}, \cite[Lemma~3.26]{Wu26}}]\label{prop:wu-main}
\lean{EnvelopeProperty.exists_antitone_fs_growth_of_local}
\leanok
\uses{conv:wu-setting, lem:wu-collar, thm:wu-pl-approx, lem:wu-growth-diagonal, lem:fs-basic}
Let $v\colon Y^\an\to[-\infty,\infty)$, $g\in\FS(Y)$ and $C\in\R$ with $v\le g+C$, and assume that every compact
subset of $Y^\an$ has an open neighbourhood on which $v$ is the limit of a decreasing sequence of convex PL
functions. For every rational $c>1$ there are $G_m\in\FS(Y)$ decreasing to $v$ and constants $A_m,B_m$ with
$cg-A_m\le G_m\le cg+B_m$.
\end{proposition}

\begin{proof}
Replacing $g$ by $g-b_0\ge0$ changes $cg$ by the constant $cb_0$ and $C$ by $b_0$; so assume $g\ge0$. Fix rationals
$1<c_1<c_2<c_3<c$. For $n\in\N$ let $t_n\coloneqq(C+n+2)/(c_1-1)$, choose an open neighbourhood $U_n$ of the
compact set $\{g\le t_n+1\}$ and convex PL functions $v_{n,j}$ on $U_n$ decreasing to $v$ there. Lemma~\ref{lem:wu-collar},
with $c_1$, $C_1=-n$ and $t=t_n$, gives (after discarding finitely many indices) functions $W_{n,j}$, standard at
infinity with respect to $h_n\coloneqq c_1g-n\in\FS(Y)$, decreasing in $j$ to $w_n\coloneqq\max\{v,c_1g-n\}$.
By Theorem~\ref{thm:wu-pl-approx} with the rational $c_2/c_1>1$, each $W_{n,j}$ is the limit of a decreasing
sequence $(F_{n,j,i})_i$ in $\FS(Y)$ with $F_{n,j,0}\le(c_2/c_1)h_n+C_{n,j}\le c_2g+C_{n,j}$. For fixed $n$,
Lemma~\ref{lem:wu-growth-diagonal} (with $c_2<c_3$, applied to the sequence $(W_{n,j})_j$ and the approximants
$(F_{n,j,i})_i$) shows that $w_n$ is the limit of a decreasing sequence $(G_{n,m})_m$ in $\FS(Y)$ with
$G_{n,0}\le c_3g+B_{n}$. The $w_n$ decrease to $v$, because $c_1g-n$ decreases to $-\infty$. A second application
of Lemma~\ref{lem:wu-growth-diagonal} (with $c_3<c$, to $(w_n)$ and $(G_{n,m})_m$) gives the assertion.
\end{proof}

\begin{theorem}[Growth-controlled approximation {\cite[Lemma~5.11]{FHX}, \cite[Lemma~3.25]{Wu26}}]\label{thm:wu-growth-control}
\lean{EnvelopeProperty.ConvexPLApproximation.growthControl}
\leanok
\uses{prop:wu-main, lem:fs-basic, def:fs}
Let $u\in\Rcal(Y)$ and $g\in\FS(Y)$, $C\in\R$ with $u\le g+C$. For every rational $c>1$ there are $u_j\in\FS(Y)$
decreasing to $u$ such that each $u_j-cg$ is bounded.
\end{theorem}

\begin{proof}
The functions of a decreasing sequence in $\FS(Y)$ with limit $u$ are convex PL on $Y^\an$
(Lemma~\ref{lem:fs-basic}), so the hypothesis of Proposition~\ref{prop:wu-main} holds with the neighbourhood
$Y^\an$ of every compact set; the proposition gives the $u_j$, with $u_j-cg$ bounded on both sides.
\end{proof}

\begin{theorem}[Local-to-global regularization {\cite[Lemma~5.11]{FHX}, \cite[Lemma~3.26]{Wu26}}]\label{thm:wu-local-to-global}
\lean{EnvelopeProperty.ConvexPLApproximation.localToGlobal, EnvelopeProperty.ConvexPLApproximation.toWuRegularization}
\leanok
\uses{prop:wu-main, def:fs, def:convex-pl}
Let $v\colon Y^\an\to[-\infty,\infty)$ be a function of logarithmic growth, not identically $-\infty$, such that
every compact subset of $Y^\an$ has an open neighbourhood on which $v$ is the limit of a decreasing sequence of
convex PL functions. Then $v\in\Rcal(Y)$.
\end{theorem}

\begin{proof}
Choose $g\in\FS(Y)$ and $C$ with $v\le g+C$. Proposition~\ref{prop:wu-main} with $c=2$ exhibits $v$ as the limit
of a decreasing sequence in $\FS(Y)$, and $v\not\equiv-\infty$.
\end{proof}

\begin{remark}\label{rmk:wu-affine}
\uses{thm:wu-growth-control, thm:wu-local-to-global, lem:wu-pl-special-case}
Theorems~\ref{thm:wu-growth-control} and~\ref{thm:wu-local-to-global} are the statements of
\cite[Lemma~5.11(5),(6)]{FHX}, the former with the two-sided bound of the source. Together with
Lemma~\ref{lem:wu-pl-special-case} they show that these two statements are jointly equivalent to
Theorem~\ref{thm:wu-pl-approx}. The source proves them on a projective compactification, through
\cite[Lemma~5.10]{FHX}, PSH2 \cite[Theorem~4.5]{BJ22} and countable regularization in a fixed ample class
\cite[Corollary~12.18(iii)]{BJ22}. Here the truncation and gluing of \cite[Lemma~5.10(2)]{FHX} and
\cite[Lemma~3.24(3)]{Wu26} is carried out on $Y^\an$ (Lemma~\ref{lem:wu-collar}); compactness of the
compactification is replaced by the gap between growth rates $c'<c$ in Lemma~\ref{lem:wu-growth-diagonal}, whose
truncations $\max\{F,cg-m\}$ agree with $cg-m$ outside compact sets.
\end{remark}

\subsection{A normal compactification}\label{subsec:wu-compactification}

\begin{convention}\label{conv:wu-graded}
\lean{EnvelopeProperty.IsTGraded, EnvelopeProperty.gradedPiece, EnvelopeProperty.veronese, EnvelopeProperty.homLoc, EnvelopeProperty.ProjBerk.zeroLocus, EnvelopeProperty.IsHomogeneousIdeal, EnvelopeProperty.ProjSubvariety, EnvelopeProperty.ProjSubvariety.an}
\leanok
\uses{conv:wu-setting, def:pn-an, lem:bj-chart}
For a $k$-algebra $A$, the ring $A[T]$ is graded by the degree in $T$. A \emph{graded subalgebra} $B\subseteq A[T]$
is a $k$-subalgebra containing the homogeneous components of its elements; $B_d\coloneqq B\cap AT^d$. For
homogeneous $\sigma\in B$, $B_{(\sigma)}$ denotes the degree-zero part of $B[\sigma^{-1}]$, and for $e\ge1$ the
Veronese subalgebra $B^{(e)}\coloneqq\bigoplus_dB_{de}$ is graded by $d$. For $M\in\N$ put $R_M=k[x_0,\dots,x_M]$
and use the model of $(\PP^M)^\an$ of Definition~\ref{def:pn-an}. For a homogeneous prime ideal $I\subseteq R_M$
not containing all $x_i$, let $V(I)^\an\coloneqq\{v\in(\PP^M)^\an:|f|_v=0\ \forall f\in I\}$. For a form $\sigma$
of degree $d\ge1$, $U_\sigma=\{|\sigma|\ne0\}$ and $\iota_\sigma$ are as in Lemma~\ref{lem:bj-chart}. All
elements of $R_M$ have absolute value $\le1$ at points of $(\PP^M)^\an$.
\end{convention}

\begin{lemma}[Normal form of Fubini--Study functions {\cite[Proposition~3.17]{Wu26}, \cite[Lemma~5.9]{FHX}}]\label{lem:wu-fs-normal-form}
\lean{EnvelopeProperty.FSData.NormalForm, EnvelopeProperty.FSData.nonempty_normalForm}
\leanok
\uses{conv:wu-setting, def:fs}
For every $h\in\FS(Y)$ there are $N\in\N_{>0}$, generators $y_1,\dots,y_m$ of $A$ and
$\lambda_0,\dots,\lambda_m\in\Q$ with $Nh=\max\{\lambda_0,\log|y_1|+\lambda_1,\dots,\log|y_m|+\lambda_m\}$.
\end{lemma}

\begin{proof}
Write $h=\max\{a_1\log|f_1|+b_1,\dots,a_r\log|f_r|+b_r,b_0\}$ and choose $N$ with $Na_i\in\N_{\ge1}$ for all $i$,
and $M\in\Q$ with $M\ge|Nb_i|$ for all $i\ge0$. Since $Na_i\ge1$, $\log|f_i|\le\max\{Na_i\log|f_i|,0\}$ (separate
$|f_i|\ge1$ and $|f_i|<1$), and $\max\{Na_i\log|f_i|,0\}\le\max\{Na_i\log|f_i|+Nb_i,Nb_0\}+M\le Nh+M$. Hence
$Nh=\max\{\log|f_i^{Na_i}|+Nb_i,\ \log|f_i|-M,\ Nb_0\}$, and the family $(f_i^{Na_i})_i\cup(f_i)_i$ contains the
generators among the $f_i$.
\end{proof}

\begin{lemma}[Normalized Rees algebra {\cite[Lemma~5.9]{FHX}}]\label{lem:wu-graded-closure}
\lean{EnvelopeProperty.normalizedRees_spec, EnvelopeProperty.isGraded_normalizedRees, EnvelopeProperty.fg_normalizedRees, EnvelopeProperty.finiteDimensional_gradedPiece_normalizedRees, EnvelopeProperty.isIntegrallyClosed_normalizedRees, EnvelopeProperty.exists_algebraMap_eq_of_C_mem_normalizedRees, EnvelopeProperty.homLocNormalizedReesEquiv, EnvelopeProperty.exists_pow_eq_sum_reesGen}
\leanok
\uses{conv:wu-graded, thm:bj-normalization-finite}
Let $y_1,\dots,y_m$ generate $A$, and set
\[
S\coloneqq k[T,y_1T,\dots,y_mT]\subseteq A[T].
\]
Let $S'$ be the integral closure of $S$ in $A[T]$. Then:
\begin{enumerate}
\item $S'$ is a graded subalgebra of $A[T]$;
\item $S'$ is a finitely generated $k$-algebra, and each $S'_d$ is finite dimensional;
\item $S'$ is an integrally closed domain;
\item $S'_0=k$;
\item the map $S'_{(T)}\to A$, $bT^d/T^d\mapsto b$, is an isomorphism of $k$-algebras;
\item every homogeneous $b\in S'$ of positive degree has a power in the ideal of $S'$ generated by
$T,y_1T,\dots,y_mT$.
\end{enumerate}
\end{lemma}

\begin{proof}
(1) For $\lambda\in k^\times$ let $\theta_\lambda$ be the $A$-algebra automorphism of $A[T]$ with
$T\mapsto\lambda T$. It preserves $S$, hence $S'$. If $b=\sum_{d=0}^Db_dT^d\in S'$, then
$\theta_\lambda(b)=\sum_d\lambda^db_dT^d\in S'$ for $D+1$ distinct $\lambda\in k^\times$ ($k$ is infinite), and
inverting the Vandermonde matrix expresses each $b_dT^d$ as a $k$-linear combination of these elements.

(2) $S$ is a finitely generated $k$-domain and $k$ is perfect, so the integral closure $\widetilde S$ of $S$ in
$\operatorname{Frac}S$ is a finite $S$-module (Theorem~\ref{thm:bj-normalization-finite}). Since $y_l=y_lT/T$,
$\operatorname{Frac}S=\operatorname{Frac}(A)(T)\supseteq A[T]$, so $S'=\widetilde S\cap A[T]$ is an $S$-submodule
of a finite module over the Noetherian ring $S$, hence finite. Thus $S'$ is a finitely generated $k$-algebra,
and $S'_d$ is spanned by the products of the homogeneous components of finitely many module generators with
monomials in $T,y_lT$ of complementary degree.

(3) $A[T,T^{-1}]$ is an integrally closed domain, being a localization of the polynomial ring over the normal
domain $A$. If $z\in\operatorname{Frac}S'=\operatorname{Frac}A[T,T^{-1}]$ is integral over $S'$, it is integral
over $S$, so $z\in A[T,T^{-1}]$. The argument of (1) in $A[T,T^{-1}]$ shows that every homogeneous component
$z_dT^d$ of $z$ is integral over $S$. For $d<0$, taking degree-$nd$ components of an equation
$(z_dT^d)^n+s_1(z_dT^d)^{n-1}+\dots+s_n=0$ with $s_i\in S$ we may assume $s_i\in S_{id}=0$, so $z_d=0$. Hence
$z\in A[T]$ and $z\in S'$.

(4) For $a\in S'_0\subseteq A$, taking degree-zero components of an integral equation over $S$ shows that $a$ is
integral over $S_0=k$; as $k$ is algebraically closed and $A$ is a domain containing $k$, $a\in k$.

(5) The map is well defined and injective because $S'[T^{-1}]\subseteq A[T,T^{-1}]$. Its image is a
$k$-subalgebra containing $y_l=(y_lT)/T$, hence equal to $A$.

(6) Let $b\in S'_d$, $d>0$. Taking homogeneous components of degree $nd$ in an integral equation of $b$ over $S$
we get $b^n+s_1b^{n-1}+\dots+s_n=0$ with $s_i\in S_{id}$, which lies in the ideal of $S$ generated by
$T,y_1T,\dots,y_mT$ since $id>0$. Hence $b^n$ lies in the ideal they generate in $S'$.
\end{proof}

\begin{lemma}[Veronese subalgebras]\label{lem:wu-veronese}
\lean{EnvelopeProperty.exists_veronese_generated, EnvelopeProperty.isIntegrallyClosed_veronese, EnvelopeProperty.homLoc_veronese, EnvelopeProperty.exists_veronese_normalizedRees}
\leanok
\uses{conv:wu-graded}
Let $B\subseteq A[T]$ be a finitely generated graded subalgebra with $B_0=k$. There is $e\ge1$ with
$B_{ne}=(B_e)^n$ (the span of $n$-fold products) for all $n\ge1$; thus $B^{(e)}$ is generated by $B_e$. If $B$ is
an integrally closed domain, so is $B^{(e)}$. For homogeneous $\sigma\in B_e$, $B^{(e)}_{(\sigma)}=B_{(\sigma)}$.
\end{lemma}

\begin{proof}
Let $b_1,\dots,b_r$ be homogeneous generators of positive degrees $d_i$ (possible as $B_0=k$), let $\delta$ be the
least common multiple of the $d_i$, $\beta_i\coloneqq b_i^{\delta/d_i}\in B_\delta$ and $e\coloneqq\max(r,1)\delta$ (so $e\ge1$ also when $B=k$). Every
monomial $b^a$ equals $\mu\beta^\gamma$ with $\mu=\prod_ib_i^{\rho_i}$, $0\le\rho_i<\delta/d_i$ (divide $a_i$ by
$\delta/d_i$), and $\deg\mu=\sum_i\rho_id_i<r\delta\le e$. If $\deg b^a=ne$, then $\delta$ divides $\deg\mu$ and
$\beta^\gamma$ is a product of $(ne-\deg\mu)/\delta=(e-\deg\mu)/\delta+(n-1)e/\delta$ factors $\beta_i$. Grouping
$\mu$ with $(e-\deg\mu)/\delta$ of them and the others in $n-1$ groups of $e/\delta$ writes $b^a$ as a product of
$n$ elements of $B_e$. Since $B_{ne}$ is spanned by monomials of degree $ne$, $B_{ne}=(B_e)^n$.

Let $B$ be an integrally closed domain and $z\in\operatorname{Frac}B^{(e)}$ integral over $B^{(e)}$. Then $z\in B$.
Since $B^{(e)}\subseteq A[T^e]$, we can write $z=p(T^e)/q(T^e)$ with $p,q\in\operatorname{Frac}(A)[X]$, $q\ne0$. For
a residue class $\varrho$ modulo $e$ let $z^{(\varrho)}$ be the sum of the components of $z$ of degree $\equiv\varrho$.
Comparing components in $q(T^e)z=p(T^e)$ gives $q(T^e)z^{(\varrho)}=0$ for $\varrho\ne0$, hence $z^{(\varrho)}=0$.
So all components of $z$ have degree divisible by $e$; they lie in $B$ by gradedness, and $z\in B^{(e)}$.
Finally an element of $B_{(\sigma)}$ is $b/\sigma^j$ with $b\in B_{je}$, which lies in $B^{(e)}_{(\sigma)}$.
\end{proof}

\begin{definition}[The compactification attached to a Fubini--Study function {\cite[Lemma~5.9]{FHX}}]\label{def:wu-compactification}
\lean{EnvelopeProperty.FSData.Compactification, EnvelopeProperty.FSData.Compactification.φ, EnvelopeProperty.FSData.Compactification.proj, EnvelopeProperty.FSData.Compactification.bdry, EnvelopeProperty.FSData.Compactification.ψ, EnvelopeProperty.FSData.nonempty_compactification, EnvelopeProperty.FSData.Compactification.isIntegrallyClosed_quotient}
\leanok
\uses{lem:wu-fs-normal-form, lem:wu-graded-closure, lem:wu-veronese, conv:wu-graded}
Let $h\in\FS(Y)$. Choose $N$, $y_l$ and $\lambda_l$ by Lemma~\ref{lem:wu-fs-normal-form}, let $S\subseteq S'$ be as in
Lemma~\ref{lem:wu-graded-closure}, choose $e$ by Lemma~\ref{lem:wu-veronese} for $B=S'$, and put
$S''\coloneqq S'^{(e)}$. Let $z_0=T^e,z_1,\dots,z_M$ be a $k$-basis of $S''_1=S'_e$, and let $I$ be the kernel of
the surjective graded homomorphism $R_M\to S''$, $x_i\mapsto z_i$. Let $\rho_l\in R_M$ be linear forms mapping to
$(y_lT)^e\in S''_1$. Put $P^\an\coloneqq V(I)^\an$ and
\[
\psi\coloneqq\max\bigl\{\log|x_0|,\ \log|\rho_1|+e(\lambda_1-\lambda_0),\dots,\log|\rho_m|+e(\lambda_m-\lambda_0)\bigr\}
\quad\text{on }P^\an.
\]
\end{definition}

\begin{lemma}[Charts of the compactification]\label{lem:wu-proj-chart}
\lean{EnvelopeProperty.FSData.Compactification.isReesPresentation, EnvelopeProperty.FSData.Compactification.chart_triv, EnvelopeProperty.ProjSubvariety.isCompact_an, EnvelopeProperty.ProjBerk.mem_chartIdeal_iff, EnvelopeProperty.ProjBerk.quotChartHomeomorph, EnvelopeProperty.ProjSubvariety.chartHomeomorph, EnvelopeProperty.ProjSubvariety.logAbs_chartMap_quotChartMk, EnvelopeProperty.ProjSubvariety.sup_logAbs_chartMap_quotChartGen, EnvelopeProperty.ProjSubvariety.affineChart, EnvelopeProperty.ProjSubvariety.affineChart_triv, EnvelopeProperty.ProjBerk.IsReesPresentation}
\leanok
\uses{def:wu-compactification, conv:wu-graded, lem:bj-chart, def:pn-an, lem:wu-graded-closure, lem:wu-veronese}
Let $I\subseteq R_M$ be a homogeneous prime not containing all $x_i$, and $S''=R_M/I$.
\begin{enumerate}
\item $V(I)^\an$ is compact.
\item For a form $\sigma$ of degree $d\ge1$ with $\sigma\notin I$, $\iota_\sigma$ restricts to a homeomorphism of
$U_\sigma\cap V(I)^\an$ onto $(\Spec S''_{(\bar\sigma)})^\an$, where $\bar\sigma$ is the image of $\sigma$.
\item For a form $\tau$ of degree $nd$, $\log|\tau|-n\log|\sigma|=\log|\bar\tau/\bar\sigma^n|\circ\iota_\sigma$ on
$U_\sigma\cap V(I)^\an$, and $-\frac1d\log|\sigma|=\frac1d\max_{|\alpha|=d}\log|\bar x^\alpha/\bar\sigma|\circ\iota_\sigma$.
\item In the situation of Definition~\ref{def:wu-compactification}, $\iota_{x_0}$ identifies $U_{x_0}\cap P^\an$
with $Y^\an$, through $S''_{(z_0)}=S'_{(T)}\cong A$, and maps the trivial valuation of $S''$ to $\vtriv$.
\end{enumerate}
In Lean, (2) holds for every homogeneous $I$ (primality and $\sigma\notin I$ are used only for the description of
the kernel), and (4) is stated for any graded presentation $R_M\to A[T]$, $x_{j_0}\mapsto T^e$, with kernel $I$ and
degree-zero part $A$ (\texttt{ProjBerk.IsReesPresentation}); Definition~\ref{def:wu-compactification} provides one.
\end{lemma}

\begin{proof}
(1) $V(I)^\an$ is closed in the compact space $(\PP^M)^\an$. (2) By Lemma~\ref{lem:bj-chart}(2), $\iota_\sigma$ is a
homeomorphism of $U_\sigma$ onto $(\Spec(R_M)_{(\sigma)})^\an$ with $|\tau/\sigma^n|_{\iota_\sigma(v)}=|\tau|_v/|\sigma|_v^n$.
Since $I$ is prime and $\sigma\notin I$, the kernel of $(R_M)_{(\sigma)}\to S''_{(\bar\sigma)}$ is
$I_{(\sigma)}=\{\tau/\sigma^n:\tau\in I\text{ homogeneous of degree }nd\}$. A graded $v\in U_\sigma$ vanishes on the
homogeneous ideal $I$ if and only if it vanishes on its homogeneous elements; for such $\tau$ of degree $D$ we have
$\tau^d/\sigma^D\in I_{(\sigma)}$ and $|\tau^d/\sigma^D|_{\iota_\sigma(v)}=|\tau|_v^d/|\sigma|_v^D$. Hence $v$ vanishes on $I$
if and only if $\iota_\sigma(v)$ vanishes on $I_{(\sigma)}$. Seminorms vanishing on an ideal are the seminorms of the
quotient, with the induced topology. (3) is Lemma~\ref{lem:bj-chart}(3). (4) Combine (2) with
Lemma~\ref{lem:wu-veronese} and Lemma~\ref{lem:wu-graded-closure}(5); the trivial valuation has $|x_0|=1$, and
Lemma~\ref{lem:bj-chart}(4) applies.
\end{proof}

\begin{definition}[Fubini--Study and locally convex functions on the compactification]\label{def:wu-proj-functions}
\lean{EnvelopeProperty.ProjSubvariety.IsFiniteFS, EnvelopeProperty.ProjSubvariety.IsFiniteFS.continuous, EnvelopeProperty.ProjSubvariety.IsLocallyConvexAt, EnvelopeProperty.ProjSubvariety.IsLocallyConvex, EnvelopeProperty.ProjSubvariety.IsFiniteFS.isLocallyConvex, EnvelopeProperty.ProjSubvariety.IsLocallyConvex.isLocallyConvexAll, EnvelopeProperty.ProjSubvariety.isLocallyConvexAll_iff, EnvelopeProperty.ProjSubvariety.isConvexPLOn_comp_chartInv}
\leanok
\uses{conv:wu-graded, lem:wu-proj-chart, def:convex-pl}
Let $I\subseteq R_M$ be as in Lemma~\ref{lem:wu-proj-chart}, $S''=R_M/I$ and $P^\an=V(I)^\an$.
\begin{enumerate}
\item A \emph{finite Fubini--Study function of class $\mathcal O(1)$} on $P^\an$ is a function
$\Psi=\frac1m\max_\ell\{\log|\tau_\ell|+\mu_\ell\}$ with $m\ge1$, finitely many forms $\tau_\ell$ of degree $m$ and
$\mu_\ell\in\Q$, such that at every point of $P^\an$ some $\tau_\ell$ does not vanish. It is continuous and real valued.
\item Let $a\in\Q_{>0}$. A function $\varphi\colon P^\an\to\R$ is \emph{locally $\mathcal O(a)$-convex} if every
$x\in P^\an$ has an open neighbourhood $W$ such that for some form $\sigma\notin I$ of degree $d\ge1$ with
$W\subseteq U_\sigma$, the function $\bigl(\varphi-\frac ad\log|\sigma|\bigr)\circ\iota_\sigma^{-1}$ is convex PL on
$\iota_\sigma(W)\subseteq(\Spec S''_{(\bar\sigma)})^\an$.
\end{enumerate}
By Lemma~\ref{lem:wu-proj-chart}(3) and the stability of convex PL functions under adding rational multiples of
$\log$ of invertible functions (Lemma~\ref{lem:local-basic}), (2) is the notion of locally $\mathcal O(a)$-convex
function of \cite[Definition~5.3]{FHX}, in the form given by the remark following it.
\end{definition}

\begin{lemma}[Boundary presentation {\cite[Lemma~5.9]{FHX}, \cite[Proposition~3.17]{Wu26}}]\label{lem:wu-boundary-presentation}
\lean{EnvelopeProperty.FSData.Compactification.isFiniteFS_ψ, EnvelopeProperty.FSData.Compactification.exists_val_bdry_ne_zero, EnvelopeProperty.FSData.Compactification.eval_chart, EnvelopeProperty.FSData.Compactification.ψ_eq, EnvelopeProperty.FSData.exists_boundaryPresentation}
\leanok
\uses{lem:wu-fs-normal-form, def:wu-compactification, lem:wu-proj-chart, lem:wu-graded-closure, def:wu-proj-functions}
In the situation of Definition~\ref{def:wu-compactification}, $\psi$ is a finite Fubini--Study function of class
$\mathcal O(1)$ on $P^\an$, and on $U_{x_0}\cap P^\an=Y^\an$
\[
h=\frac1{eN}\bigl(\psi-\log|x_0|\bigr)+\frac{\lambda_0}N .
\]
\end{lemma}

\begin{proof}
Let $g_0=T$ and $g_l=y_lT$, and let $J$ (resp.\ $J'$) be the ideal of $S'$ generated by the $g_l$ (resp.\ the
$g_l^e$). A product of $p\coloneqq(m+1)(e-1)+1$ of the $g_l$ contains some $g_l$ at least $e$ times, so $J^p\subseteq J'$.
Let $b\in S''_1=S'_e$. By Lemma~\ref{lem:wu-graded-closure}(6), $b^n\in J$ for some $n$, so $b^{np}\in J'$; keeping
the components of degree divisible by $e$ of the coefficients, $b^{np}=\sum_lg_l^ec_l$ with $c_l\in S''$. If
$v\in P^\an$ vanished at $x_0$ and at all $\rho_l$, then $|b|_v^{np}\le\max_l|g_l^e|_v|c_l|_v=0$ for all
$b\in S''_1$, contradicting $\max_i|x_i|_v=1$. So $\psi$ is finite.

On $U_{x_0}$, Lemma~\ref{lem:wu-proj-chart}(3),(4) gives $\log|\rho_l|-\log|x_0|=\log|(y_lT)^e/T^e|=e\log|y_l|$. Hence
$\psi-\log|x_0|=\max\{0,\max_l(e\log|y_l|+e\lambda_l-e\lambda_0)\}=e(Nh-\lambda_0)$ by
Lemma~\ref{lem:wu-fs-normal-form}.
\end{proof}

\begin{lemma}[Restriction of Fubini--Study functions {\cite[Lemma~5.10]{FHX}, \cite[Proposition~3.17]{Wu26}}]\label{lem:wu-fs-restriction}
\lean{EnvelopeProperty.ProjSubvariety.exists_fsData_eval_affineChart, EnvelopeProperty.ProjSubvariety.adjoin_range_coeff_map_X}
\leanok
\uses{def:wu-compactification, lem:wu-proj-chart, def:wu-proj-functions, def:fs}
In the situation of Definition~\ref{def:wu-compactification}, let $\Psi$ be a finite Fubini--Study function of class
$\mathcal O(1)$ on $P^\an$. Then $\Psi-\log|x_0|$, as a function on $Y^\an$, lies in $\FS(Y)$.
\end{lemma}

\begin{proof}
Write $\Psi=\frac1m\max_\ell\{\log|\tau_\ell|+\mu_\ell\}$ and let $f_\ell\in A$ correspond to $\bar\tau_\ell/z_0^m$; then
$\Psi-\log|x_0|=\frac1m\max_\ell\{\log|f_\ell|+\mu_\ell\}$ on $Y^\an$ (Lemma~\ref{lem:wu-proj-chart}(3)). Let $u_i\in A$
correspond to $z_i/z_0$; they generate $A$, because $S''_{(z_0)}$ consists of the $b/z_0^j$ with $b\in S''_j=(S''_1)^j$
(Lemma~\ref{lem:wu-veronese}). As $\Psi$ is continuous and real valued on the compact space $P^\an$, there is $C_0\in\Q$
with $\Psi\ge-C_0$. Since $\log|x_i|\le0$ and $-\log|x_0|\ge0$, we have
$\log|u_i|-C_0=\log|x_i|-\log|x_0|-C_0\le\Psi-\log|x_0|$ and $-C_0\le\Psi-\log|x_0|$. Therefore
\[
\Psi-\log|x_0|=\max\Bigl\{\tfrac1m\max_\ell(\log|f_\ell|+\mu_\ell),\ \max_i(\log|u_i|-C_0),\ -C_0\Bigr\}\in\FS(Y).
\]
\end{proof}

\begin{theorem}[Uniform approximation of convex PL functions standard at infinity {\cite[Lemma~3.24]{Wu26}, \cite[Lemma~5.10]{FHX}}]\label{thm:wu-pl-approx-uniform}
\lean{EnvelopeProperty.LocalPositivity.exists_fsData_abs_sub_le}
\leanok
\uses{def:wu-standard, def:wu-compactification, lem:wu-boundary-presentation, lem:wu-proj-chart, def:wu-proj-functions, thm:wu-pos, lem:wu-fs-restriction, lem:local-basic, lem:fs-basic}
Let $w$ be standard at infinity with respect to $h\in\FS(Y)$. For every $\epsilon>0$ there is $v\in\FS(Y)$ with
$|v-w|\le\epsilon$ on $Y^\an$.
\end{theorem}

\begin{proof}
Take the data of Definition~\ref{def:wu-compactification} for $h$, identify $Y^\an$ with $U_{x_0}\cap P^\an$
(Lemma~\ref{lem:wu-proj-chart}(4)), and let $E\subseteq Y^\an$ be compact with $w=h$ off $E$; $E$ is closed in
$P^\an$. Define $\varphi\colon P^\an\to\R$ by $\varphi\coloneqq eNw-e\lambda_0+\log|x_0|$ on $U_{x_0}\cap P^\an$ and
$\varphi\coloneqq\psi$ on $P^\an\setminus E$. On the overlap $w=h$, so both definitions agree by
Lemma~\ref{lem:wu-boundary-presentation}. Both sets are open and cover $P^\an$, and both expressions are continuous
and real valued there, so $\varphi$ is continuous.

$\varphi$ is locally $\mathcal O(1)$-convex. Near a point of $U_{x_0}$ take $\sigma=x_0$: the function
$(\varphi-\log|x_0|)\circ\iota_{x_0}^{-1}=eNw-e\lambda_0$ is convex PL on $Y^\an$ (Lemma~\ref{lem:local-basic}). Near a
point $x\notin E$ take $\sigma\in\{x_0,\rho_1,\dots,\rho_m\}$ with $|\sigma|_x\ne0$ (Lemma~\ref{lem:wu-boundary-presentation})
and $W\coloneqq U_\sigma\setminus E$. By Lemma~\ref{lem:wu-proj-chart}(3), $(\psi-\log|\sigma|)\circ\iota_\sigma^{-1}$ is the
maximum over $\rho\in\{x_0,\rho_l\}$ of $\log|\bar\rho/\bar\sigma|$ plus rational constants, and the term $\rho=\sigma$ is
constant. Near each point the terms with $|\bar\rho/\bar\sigma|=0$ there are smaller than this constant and may be omitted,
and the remaining $\bar\rho/\bar\sigma$ do not vanish nearby; so the function is convex PL on $\iota_\sigma(W)$, as in
Lemma~\ref{lem:fs-basic}.

By Theorem~\ref{thm:wu-pos} there is a finite Fubini--Study function $\Psi$ of class $\mathcal O(1)$ with
$|\varphi-\Psi|\le eN\epsilon$ on $P^\an$. By Lemma~\ref{lem:wu-fs-restriction}, $\Psi-\log|x_0|\in\FS(Y)$, hence
$v\coloneqq(\Psi-\log|x_0|+e\lambda_0)/(eN)\in\FS(Y)$, and on $Y^\an$ we get $|v-w|=|\Psi-\varphi|/(eN)\le\epsilon$.
\end{proof}

\begin{remark}\label{rmk:wu-compactification}
\uses{thm:wu-pl-approx-uniform, def:wu-compactification}
$\operatorname{Proj}S''$ is the normalization of the closure of $Y$ in $\PP^m$ under the embedding by
$(y_1,\dots,y_m)$, re-embedded by $S''_1$; this is the compactification of \cite[Lemma~5.9]{FHX}, built from
\cite[Proposition~3.17(2)]{Wu26}, with the normalization realized as a graded integral closure. Normality is needed
for positivity (\S\ref{subsec:wu-positivity}), whose proof describes PL functions by test configurations of a normal
variety. The functions $\varphi$ fed to positivity are continuous and equal to a finite Fubini--Study function near
$\{|x_0|=0\}$, so the decreasing-limit and countable-regularization parts of \cite[Lemma~5.10]{FHX} do not occur:
Lemma~\ref{lem:wu-fs-restriction} is \cite[Lemma~5.10(1)]{FHX} for finite Fubini--Study functions, and the truncation of
\cite[Lemma~5.10(2)]{FHX} is Lemma~\ref{lem:wu-collar}.
\end{remark}

\subsection{Positivity}\label{subsec:wu-positivity}

\begin{convention}\label{conv:wu-bj-global}
\notready
\uses{def:wu-proj-functions, lem:wu-proj-chart}
In this subsection $I\subseteq R_M$ is a homogeneous prime not containing all $x_i$ such that $S''=R_M/I$ is an
integrally closed domain with $S''_0=k$; $P=\operatorname{Proj}S''$ is a normal projective variety and
$L=\mathcal O_P(1)$ is very ample. As for Definition~\ref{def:pn-an}, $V(I)^\an$ is identified with the Berkovich
analytification $P^\an$. Following the author's decision D4, $\PSH(P,L)$ is read as Definition~\ref{def:pg-subvariety-psh}
(classes in $\R\cdot\mathcal O(1)$; it agrees with \cite[Definition~4.1]{BJ22} by openness of the ample cone,
Remark~\ref{rmk:pg-subvariety-psh-comparison}, which is not formalized). We use the notions of \cite{BJ22} recalled in \cite[\S2]{FHX}: the $\Q$-vector space
$\operatorname{PL}(P^\an)$ of differences of finite Fubini--Study functions of class $\mathcal O(a)$, $a\in\Q_{>0}$;
test configurations $(\mathcal X,\mathcal L)$ for $(P,rL)$ and their PL functions $\phi_{\mathcal L}$; and the class
$\PSH(P,L)$ of \cite[Definition~4.1]{BJ22}. Since $S''$ is normal, its forms of degree $m$ are the sections of
$\mathcal O_P(m)$, so the finite Fubini--Study functions of class $\mathcal O(1)$ are the functions of
$\mathcal H(L)$ in the notation of \cite{BJ22}.
\end{convention}

\begin{lemma}[Density of PL functions {\cite[Theorem~2.2]{BJ22}}]\label{lem:wu-pl-dense}
\notready
\uses{conv:wu-bj-global, thm:bj-pl-dense, lem:bj-fs-basic}
$\operatorname{PL}(P^\an)$ contains the restrictions of the rational PL functions on $(\PP^M)^\an$, is stable under
finite maxima and minima, and is dense in $C^0(P^\an)$ for uniform convergence.
\end{lemma}

\begin{proof}
A finite Fubini--Study function on $(\PP^M)^\an$ restricts to a finite one on $P^\an$, so restrictions of rational
PL functions are PL. Finite Fubini--Study functions are stable under sums and finite maxima after raising to a common
class (Lemma~\ref{lem:bj-fs-basic}), and
$\max\{\varphi_1-\psi_1,\varphi_2-\psi_2\}=\max\{\varphi_1+\psi_2,\varphi_2+\psi_1\}-(\psi_1+\psi_2)$; minima are
$-\max\{-a,-b\}$. Given $f\in C^0(P^\an)$, extend it to $F\in C^0((\PP^M)^\an)$ by Tietze's theorem
($(\PP^M)^\an$ is compact Hausdorff and $P^\an$ is closed), approximate $F$ uniformly by a rational PL function
(Theorem~\ref{thm:bj-pl-dense}) and restrict.
\end{proof}

\begin{lemma}[Locally convex functions are PL {\cite[Lemma~5.5]{FHX}}]\label{lem:wu-locally-convex-pl}
\lean{EnvelopeProperty.ProjSubvariety.IsLocallyConvex.isPL}
\leanok
\uses{conv:wu-bj-global,def:wu-proj-functions,lem:lp-pl-sheaf,lem:lp-chart-log-germ,lem:lp-log-germ}
Every locally $\mathcal O(a)$-convex function $\varphi\colon P^\an\to\R$ belongs to $\operatorname{PL}(P^\an)$.
\end{lemma}

\begin{proof}
At a point $x$, choose the finite tropical-affine presentation of local convexity in a chart
$U_\sigma$. A homogeneous form $s$ nonzero at $x$ has a global PL logarithmic germ: choose
$c\in\Q$ below $\log|s(x)|$ and use the finite rational Fubini--Study function
$\max(\log|s|,c)$. This also handles homogeneous forms of degree zero.
Write each regular chart function as $f_i=\overline\tau_i/\overline\sigma^{N_i}$.
Since it is nonzero at $x$, both numerator and denominator logarithms have global PL germs;
their difference gives a global PL representative of $\log|f_i|$ near $x$.
On the common neighbourhood of these finitely many germs, $\varphi$ is a finite maximum of
rational linear combinations of them, plus $(a/\deg\sigma)\log|\sigma|$.
Thus $\varphi$ is locally equal to global PL functions. Lemma~\ref{lem:lp-pl-sheaf} glues them
exactly. This proves PL membership; the positivity criterion below remains a separate step.
\end{proof}

\begin{theorem}[Imported: PL functions and test configurations {\cite[Theorem~2.7]{BJ22}, \cite[\S2.3]{FHX}}]\label{thm:wu-bj-testconfig}
\notready
\uses{conv:wu-bj-global}
For every $\varphi\in\operatorname{PL}(P^\an)$ there are an integer $r\ge1$ and a test configuration
$(\mathcal X,\mathcal L)$ for $(P,rL)$, with $\mathcal L$ a line bundle and $\mathcal X$ normal, dominating
$P\times\A^1$ through a projective birational morphism that is an isomorphism outside the central fibre, such that
$r\varphi=\phi_{\mathcal L}$.
\end{theorem}

\begin{theorem}[Imported: nef test configurations give psh functions {\cite[Definition~3.2, Lemma~3.1, Lemma~4.3, Lemma~3.7]{BJ22}}]\label{thm:wu-bj-nef-psh}
\notready
\uses{conv:wu-bj-global, thm:wu-bj-testconfig}
In the situation of Theorem~\ref{thm:wu-bj-testconfig}, if $\mathcal L$ has nonnegative degree on every integral curve
contained in the central fibre $\mathcal X_0$, then $\varphi\in\PSH(P,L)$.
\end{theorem}

\begin{theorem}[Imported: positivity criterion of Chambert-Loir--Ducros {\cite[Lemma~5.5]{FHX}}]\label{thm:wu-cld-nef}
\notready
\uses{conv:wu-bj-global, thm:wu-bj-testconfig, def:wu-proj-functions}
In the situation of Theorem~\ref{thm:wu-bj-testconfig}, if $r\varphi$ is locally $\mathcal O(r)$-convex, then
$\mathcal L$ has nonnegative degree on every integral curve contained in $\mathcal X_0$.
\end{theorem}

\begin{theorem}[Uniform approximation of continuous psh functions {\cite[Theorem~4.15]{BJ22}, \cite[Proposition~2.11]{Wu26}}]\label{thm:wu-bj-uniform}
\lean{EnvelopeProperty.ProjSubvariety.IsPSH.exists_isFiniteFS_one_abs_sub_le, EnvelopeProperty.ProjSubvariety.IsPSH.exists_seq_isFiniteFS_one_abs_sub_le}
\leanok
\uses{conv:wu-bj-global, def:wu-proj-functions, lem:pg-uniform-approximation}
If $\varphi\colon P^\an\to\R$ is continuous and $\varphi\in\PSH(P,L)$, then for every $\epsilon>0$ there is a finite
Fubini--Study function $\Psi$ of class $\mathcal O(1)$ with $|\varphi-\Psi|\le\epsilon$ on $P^\an$.
\end{theorem}

\begin{proof}
With $\PSH(P,L)$ read as in Definition~\ref{def:pg-subvariety-psh} (author's decision D4), this is
Lemma~\ref{lem:pg-uniform-approximation}; it is a theorem, not an import.
\end{proof}

\begin{lemma}[Locally convex functions are psh {\cite[Lemma~5.5]{FHX}}]\label{lem:wu-local-positivity}
\lean{EnvelopeProperty.LocalPositivity, EnvelopeProperty.LocalPositivity.instOfIsAlgClosed}
\leanok
\uses{thm:lpn-local-positivity-instance, def:wu-proj-functions, def:pg-subvariety-psh}
Let $k$ be algebraically closed. Every locally $\mathcal O(1)$-convex function $\varphi\colon P^\an\to\R$ belongs to
$\PSH(P,L)$. In Lean this statement (for every projectively normal $V(I)$ and every class $\mathcal O(a)$, $a\in\Q_{>0}$,
with local convexity in the form of \cite[Definition~5.3]{FHX}: for every admissible pair) is the class
\texttt{LocalPositivity}~$k$; it holds for every algebraically closed field $k$.
\end{lemma}

\begin{proof}
This is Theorem~\ref{thm:lpn-local-positivity-instance}: every locally convex function admits a native potential
family (Theorem~\ref{thm:lpn-native-potential-family-exists}, through the normalized test configuration of its
piecewise linear model, nefness of the central fibre and the exact finite Fubini--Study functions of the model
potential), and functions with a native potential family are psh (Theorem~\ref{thm:lpn-local-positivity}).
\end{proof}

\begin{theorem}[Positivity, uniform form {\cite[Lemma~5.5]{FHX}, \cite[Theorem~3.23]{Wu26}, \cite[Theorem~4.15]{BJ22}}]\label{thm:wu-pos}
\lean{EnvelopeProperty.LocalPositivity.exists_isFiniteFS_abs_sub_le, EnvelopeProperty.LocalPositivity.isPSH_of_isLocallyConvex, EnvelopeProperty.ProjSubvariety.IsLocallyConvex.continuous}
\leanok
\uses{lem:wu-local-positivity, thm:wu-bj-uniform, lem:pg-uniform-approximation, conv:wu-bj-global, def:wu-proj-functions}
Let $I\subseteq R_M$ be as in Convention~\ref{conv:wu-bj-global}. For every continuous locally $\mathcal O(1)$-convex
function $\varphi\colon V(I)^\an\to\R$ and every $\epsilon>0$ there is a finite Fubini--Study function $\Psi$ of class
$\mathcal O(1)$ with $|\varphi-\Psi|\le\epsilon$ on $V(I)^\an$.
\end{theorem}

\begin{proof}
$\varphi\in\PSH(P,L)$ by Lemma~\ref{lem:wu-local-positivity}; apply Theorem~\ref{thm:wu-bj-uniform}.
\end{proof}

\begin{remark}\label{rmk:wu-positivity}
\uses{thm:wu-pos, thm:wu-bj-testconfig, thm:wu-bj-nef-psh, thm:wu-cld-nef, thm:wu-bj-uniform}
The compactifications of Definition~\ref{def:wu-compactification} satisfy the hypotheses of
Convention~\ref{conv:wu-bj-global} by Lemmas~\ref{lem:wu-graded-closure}(3),(4) and~\ref{lem:wu-veronese}.
Theorem~\ref{thm:wu-bj-testconfig} is \cite[Theorem~2.31]{BJ22} with the passage to a normal model of \cite[\S2.3]{FHX};
Theorem~\ref{thm:wu-bj-nef-psh} is \cite[Definition~3.2 and Lemma~4.3]{BJ22}, whose proof uses relative ampleness on test
configurations and the openness of the ample cone; Theorem~\ref{thm:wu-bj-uniform} is part~(iii) of the regularization
theorem of \cite[\S4]{BJ22} (arXiv label \texttt{thm:pshample}) together with the corollary following it, which uses
Dini's lemma, $\mathcal H^{\mathrm{dom}}(L)$ being contained in $\mathcal H(L)$.
Theorem~\ref{thm:wu-cld-nef} is the last paragraph of the proof of \cite[Lemma~5.5]{FHX}: after base change to
$K=k((\varpi))$, the model metric defined by $\mathcal L$ has local potentials $h\circ(\log|f_1|,\dots,\log|f_r|)$ with
$h$ convex and $f_i$ invertible, which are psh in the sense of \cite[Corollaire~18.8.9]{CLD12}; its curvature current
is positive \cite[Proposition~8.3.1]{CLD12}, the residual line bundles are nef
\cite[Th\'eor\`eme~18.8.4, Corollaire~18.8.8]{CLD12}, and $(C,\mathcal L|_C)$ is a premodel of a residual line bundle
\cite[\S17.7.9, Lemme~17.7.10]{CLD12}. Countability of $k$ is used nowhere in this chapter.
\end{remark}
